% ============================================================================
%  Capitolo 9 -- Elementi di algebra di Boole, parte 2           (Lezione 09)
% ============================================================================
\chapter{Algebra di Boole: funzioni, forme canoniche e minimizzazione}
\label{ch:boole2}

\epigraphline{Povero me! Povero me! Arriverò troppo tardi!}
  {Lewis Carroll, \emph{Alice nel Paese delle Meraviglie}, 1865}

\begin{chapgoals}
Alla fine di questo capitolo saprete
\begin{itemize}
\item definire letterale, clausola, fattore, mintermine e maxtermine;
\item dire quante funzioni booleane di $n$ variabili esistono, e nominare quelle di due variabili;
\item scrivere una funzione in forma canonica somma di prodotti e prodotto di somme a partire dalla sua tabella di verità;
\item minimizzare una funzione fino a quattro variabili con una mappa di Karnaugh;
\item ricavare il semisommatore e il sommatore completo, e spiegare come trentadue di essi eseguono \ins{add}.
\end{itemize}
\end{chapgoals}

% ----------------------------------------------------------------------------
\section{Riscaldamento sulla macchina}

Due registri contengono un bit ciascuno, 0 o 1. La loro somma può essere 0, 1
o 2, e richiede quindi due bit. Il Listato~\ref{lst:b2:warm} la calcola senza
alcuna istruzione aritmetica.

\asmfile[caption={Somma di un bit senza \ins{add}.},label=lst:b2:warm]{code/l09/s04.S}

Due istruzioni logiche sommano due bit. Il commento del listato è una tabella
di verità con due uscite, e la coppia di porte che essa descrive si chiama
\term{semisommatore}. Il Capitolo~\ref{ch:integers} affermava che a un
processore, per l'aritmetica, basta un sommatore: questo è il suo primo bit. Ma
come si trovano le porte quando è data solo la tabella? È il metodo di questo
capitolo: tabella di verità, forma canonica, minimizzazione.

% ----------------------------------------------------------------------------
\section{Funzioni booleane}

\subsection{Il vocabolario delle espressioni}

\begin{definition}{Letterale, clausola, fattore}
Un \term{letterale} è una variabile booleana, o il suo complemento, così come
compare in un'espressione: $a$, $\nott{b}$, $x$. Una \term{clausola}, o
\term{termine prodotto}, è il prodotto logico di due o più letterali:
$a\,\nott{b}\,x$. Un \term{fattore}, o \term{termine somma}, è la somma logica
di due o più letterali: $(a+\nott{x})$, $(b+c)$.
\end{definition}

Clausole e fattori con più di un letterale sono funzioni di livello~1: uno
strato di porte. Ogni occorrenza di una variabile conta come un letterale, e il
numero dei letterali misura il costo di un'espressione.

\subsection{Quante funzioni di $n$ variabili?}

\begin{definition}{Funzione booleana}
Una \term{funzione booleana} di $n$ variabili è un'applicazione
$f:B^{n}\to B$, con $B=\{0,1\}$: associa un'uscita binaria a ogni combinazione
di $n$ ingressi binari.
\end{definition}

\begin{theorem}{Numero delle funzioni booleane}
Esistono esattamente $2^{2^{n}}$ funzioni booleane distinte di $n$ variabili.
\end{theorem}

\noindent\emph{Dimostrazione.} Le combinazioni degli ingressi sono $2^{n}$, e
una funzione è determinata dalla scelta di un'uscita, 0 o 1, per ciascuna di
esse: è una stringa di $2^{n}$ bit, e di tali stringhe ce ne sono $2^{2^{n}}$.
\hfill$\square$

\smallskip
Il numero è finito, e cresce in fretta: 4 funzioni di una variabile, 16 di due,
256 di tre, 65\,536 di quattro, e più di quattro miliardi di cinque. Le
funzioni booleane sono isomorfe ai circuiti digitali: ogni circuito si può
scrivere come espressione, e ogni espressione si può costruire come circuito.

\subsection{Le sedici funzioni di due variabili}

La Tabella~\ref{tab:b2:sixteen} elenca tutte le funzioni di due variabili.
L'indice di ogni funzione è la sua colonna di uscite letta come numero binario.
Con una variabile le quattro funzioni sono la costante 0, $x$, $\nott{x}$ e la
costante 1.

\begin{table}[htb]
\centering\small
\caption{Le sedici funzioni di due variabili. Le colonne 00, 01, 10, 11 danno
l'uscita per ogni coppia di ingressi $x\,y$.}\label{tab:b2:sixteen}
\begin{tabular}{@{}lccccl@{\qquad}lccccl@{}}
\toprule
$f$ & 00 & 01 & 10 & 11 & \textbf{nome} & $f$ & 00 & 01 & 10 & 11 & \textbf{nome}\\
\midrule
$f_0$ & 0&0&0&0 & costante 0 & $f_8$ & 1&0&0&0 & \textbf{NOR}\\
$f_1$ & 0&0&0&1 & \textbf{AND}: $x\,y$ & $f_9$ & 1&0&0&1 & \textbf{XNOR} (equivalenza)\\
$f_2$ & 0&0&1&0 & $x\,\nott{y}$ (inibizione) & $f_{10}$ & 1&0&1&0 & NOT $y$\\
$f_3$ & 0&0&1&1 & $x$ (identità) & $f_{11}$ & 1&0&1&1 & $x+\nott{y}$ (implicazione $y\to x$)\\
$f_4$ & 0&1&0&0 & $\nott{x}\,y$ (inibizione) & $f_{12}$ & 1&1&0&0 & NOT $x$\\
$f_5$ & 0&1&0&1 & $y$ (identità) & $f_{13}$ & 1&1&0&1 & $\nott{x}+y$ (implicazione $x\to y$)\\
$f_6$ & 0&1&1&0 & \textbf{XOR} & $f_{14}$ & 1&1&1&0 & \textbf{NAND}\\
$f_7$ & 0&1&1&1 & \textbf{OR}: $x+y$ & $f_{15}$ & 1&1&1&1 & costante 1\\
\bottomrule
\end{tabular}
\end{table}

\subsection{Insiemi funzionalmente completi}

Qualsiasi funzione booleana si può realizzare usando solo AND, OR e NOT; per
esempio $x\oplus y=\nott{x}\,y+x\,\nott{y}$. (Il
Paragrafo~\ref{sec:b2:canonical} dimostra l'enunciato generale per
costruzione.) L'insieme $\{\text{AND},\text{OR},\text{NOT}\}$ si dice perciò
\term{funzionalmente completo}.

Per De Morgan, AND si può costruire con OR e NOT, e OR con AND e NOT: anche gli
insiemi $\{\text{AND},\text{NOT}\}$ e $\{\text{OR},\text{NOT}\}$ sono
completi. E basta una sola porta.

\begin{theorem}{Porte universali}
Il NAND da solo è funzionalmente completo, e così il NOR da solo.
\end{theorem}

\noindent\emph{Dimostrazione} per il NAND, scrivendo $x\uparrow y=\nott{xy}$:
\[
\nott{x}=x\uparrow x,\qquad
x\,y=(x\uparrow y)\uparrow(x\uparrow y),\qquad
x+y=(x\uparrow x)\uparrow(y\uparrow y),
\]
l'ultima per De Morgan. La dimostrazione per il NOR è duale. \hfill$\square$

\smallskip
Un solo tipo di porta basta a costruire qualsiasi circuito digitale, un intero
processore compreso. I chip reali sono costruiti in gran parte con NAND e NOR
perché sono le porte che richiedono meno transistor
(Capitolo~\ref{ch:electronics}).

\subsection{Tabelle di verità}

Una funzione booleana può essere definita dalla sua tabella di verità, che dà
il valore della variabile dipendente per ogni combinazione delle variabili
indipendenti. Una funzione di $n$ variabili è definita su $2^{n}$ punti, e la
sua tabella ha $2^{n}$ righe, elencate nell'ordine del conteggio: 000, 001,
010, \ldots\ La funzione $F$ della Tabella~\ref{tab:b2:F}, che vale 1 sulle
righe 3, 4, 5, 6 e 7, sarà seguita per tutto il capitolo.

\begin{table}[htb]
\centering\small
\caption{La funzione $F$, con il mintermine o il maxtermine di ogni
riga.}\label{tab:b2:F}
\begin{tabular}{@{}cccc|cll@{}}
\toprule
riga & $A$ & $B$ & $C$ & $F$ & \textbf{mintermine ($F=1$)} & \textbf{maxtermine ($F=0$)}\\
\midrule
0 & 0&0&0 & 0 & & $M_0=A+B+C$\\
1 & 0&0&1 & 0 & & $M_1=A+B+\nott{C}$\\
2 & 0&1&0 & 0 & & $M_2=A+\nott{B}+C$\\
3 & 0&1&1 & 1 & $m_3=\nott{A}\,B\,C$ & \\
4 & 1&0&0 & 1 & $m_4=A\,\nott{B}\,\nott{C}$ & \\
5 & 1&0&1 & 1 & $m_5=A\,\nott{B}\,C$ & \\
6 & 1&1&0 & 1 & $m_6=A\,B\,\nott{C}$ & \\
7 & 1&1&1 & 1 & $m_7=A\,B\,C$ & \\
\bottomrule
\end{tabular}
\end{table}

\subsection{Il semisommatore}

Il \term{semisommatore} \eng{half adder} somma due bit $A$ e $B$ e produce un
bit di somma $S$ e un riporto $C$. In binario $1+1$ dà somma 0 e riporto 1;
dalla tabella, $S$ vale 1 quando gli ingressi sono diversi e $C$ vale 1 quando
entrambi valgono 1 (Figura~\ref{fig:b2:ha}):
\[ S = A\oplus B,\qquad C=A\cdot B. \]
Il circuito è esattamente la coppia di istruzioni del riscaldamento. È «mezzo»
sommatore perché non ha un riporto in ingresso: può sommare soltanto i bit meno
significativi di due numeri.

\begin{figure}[htb]
\centering
\begin{minipage}{0.3\textwidth}\centering\small\ttfamily
\begin{tabular}{cc|cc}\textrm{$A$}&\textrm{$B$}&\textrm{$S$}&\textrm{$C$}\\\hline 0&0&0&0\\0&1&1&0\\1&0&1&0\\1&1&0&1\end{tabular}
\end{minipage}
\begin{minipage}{0.55\textwidth}\centering
\begin{circuitikz}
  \node[xor port] (x) at (3,1.2) {};
  \node[and port] (a) at (3,-0.6) {};
  \draw (x.in 1) -- ++(-1.6,0) node[left]{$A$} coordinate (A);
  \draw (x.in 2) -- ++(-1.6,0) node[left]{$B$} coordinate (B);
  \draw (x.in 1) ++(-1.0,0) node[circ]{} |- (a.in 1);
  \draw (x.in 2) ++(-0.5,0) node[circ]{} |- (a.in 2);
  \draw (x.out) -- ++(0.5,0) node[right]{$S$};
  \draw (a.out) -- ++(0.5,0) node[right]{$C$};
\end{circuitikz}
\end{minipage}
\caption{Il semisommatore.}\label{fig:b2:ha}
\end{figure}

\subsection{Il sommatore completo}

Il \term{sommatore completo} \eng{full adder} somma due bit e un riporto
entrante, $C_{in}$, e produce un bit di somma e un riporto uscente, $C_{out}$
(Figura~\ref{fig:b2:fa}). $S$ vale 1 quando un numero dispari di ingressi vale
1; $C_{out}$ vale 1 quando almeno due ingressi valgono 1:
\[ S=A\oplus B\oplus C_{in},\qquad C_{out}=A\,B+C_{in}\,(A\oplus B). \]
Il circuito è fatto di due semisommatori e un OR.

\begin{figure}[htb]
\centering
\begin{minipage}{0.32\textwidth}\centering\small\ttfamily
\begin{tabular}{ccc|cc}\textrm{$A$}&\textrm{$B$}&\textrm{$C_{in}$}&\textrm{$S$}&\textrm{$C_{out}$}\\\hline
0&0&0&0&0\\0&0&1&1&0\\0&1&0&1&0\\0&1&1&0&1\\1&0&0&1&0\\1&0&1&0&1\\1&1&0&0&1\\1&1&1&1&1\end{tabular}
\end{minipage}
\begin{minipage}{0.64\textwidth}\centering
\begin{circuitikz}[scale=0.9,transform shape]
  \node[xor port] (x1) at (2.4,2.0) {};
  \node[xor port] (x2) at (5.6,1.72) {};
  \node[and port] (a1) at (5.6,0.1) {};
  \node[and port] (a2) at (5.6,-1.5) {};
  \node[or port]  (o)  at (8.0,-0.7) {};
  \draw (x1.in 1) -- ++(-1.3,0) node[left]{$A$};
  \draw (x1.in 2) -- ++(-1.3,0) node[left]{$B$};
  \draw (x1.out) -- (x2.in 1);
  \draw (x2.in 2) -- ++(-0.6,0) coordinate (c) |- (-0.05,-2.6) node[left]{$C_{in}$};
  \draw (x1.out) ++(0.35,0) node[circ]{} |- (a1.in 1);
  \draw (c) ++(0,-1.0) node[circ]{} -- ++(0.0,0) |- (a1.in 2);
  \draw (x1.in 1) ++(-0.9,0) node[circ]{} |- (a2.in 1);
  \draw (x1.in 2) ++(-0.45,0) node[circ]{} |- (a2.in 2);
  \draw (a1.out) -| (o.in 1);
  \draw (a2.out) -| (o.in 2);
  \draw (x2.out) -- ++(2.6,0) node[right]{$S$};
  \draw (o.out) -- ++(0.3,0) node[right]{$C_{out}$};
\end{circuitikz}
\end{minipage}
\caption{Il sommatore completo: due semisommatori e un OR. I fili che si
incrociano senza un punto non sono collegati.}\label{fig:b2:fa}
\end{figure}

% ----------------------------------------------------------------------------
\section{Forme canoniche}\label{sec:b2:canonical}

Le forme canoniche sono una via sistematica dalla tabella a un'espressione.

\subsection{Forme elementari: funzioni di livello 2}

\begin{definition}{Somma di prodotti e prodotto di somme}
Un'espressione di \term{tipo P}, o \term{somma di prodotti} (SOP), è una somma
logica di clausole, come $y=a\,\nott{b}+\nott{a}\,d+b\,\nott{d}\,x$. Si
costruisce come uno strato di porte AND che alimenta una porta OR.
Un'espressione di \term{tipo S}, o \term{prodotto di somme} (POS), è un
prodotto logico di fattori, come $y=(a+\nott{b})(\nott{a}+d)(b+\nott{d}+x)$.
Si costruisce come uno strato di porte OR che alimenta una porta AND.
\end{definition}

\subsection{Mintermini e maxtermini}

\begin{definition}{Mintermine e maxtermine}
Un \term{mintermine} di $n$ variabili è un prodotto che contiene un letterale
di ogni variabile. Un \term{maxtermine} è una somma che contiene un letterale
di ogni variabile.
\end{definition}

Un mintermine vale 1 per una sola riga della tabella di verità, e un
maxtermine vale 0 per una sola riga. Il mintermine $m_i$ della riga $i$
contiene una variabile complementata dove la riga ha uno 0; il maxtermine
$M_i$ contiene una variabile complementata dove la riga ha un 1
(Tabella~\ref{tab:b2:F}). Per De Morgan, $M_i=\nott{m_i}$.

\subsection{Dalla tabella di verità alle forme canoniche}

\begin{theorem}{Forme canoniche}
Ogni funzione booleana è uguale
\begin{enumerate}[nosep]
\item alla somma dei mintermini delle righe in cui vale 1 (\term{forma canonica P});
\item al prodotto dei maxtermini delle righe in cui vale 0 (\term{forma canonica S}).
\end{enumerate}
Ciascuna delle due forme è unica, a meno dell'ordine dei termini.
\end{theorem}

\noindent\emph{Dimostrazione.} Su una riga in cui $f=1$, il mintermine di
quella riga vale 1 e così la somma; su una riga in cui $f=0$, tutti i
mintermini della somma valgono 0. Il secondo enunciato è il duale.
\hfill$\square$

\smallskip
Per la funzione $F$:
\begin{align*}
F &= \nott{A}BC+A\nott{B}\,\nott{C}+A\nott{B}C+AB\nott{C}+ABC \;=\; \sum m(3,4,5,6,7),\\
F &= (A+B+C)\,(A+B+\nott{C})\,(A+\nott{B}+C) \;=\; \prod M(0,1,2).
\end{align*}

Il teorema dimostra anche che AND, OR e NOT sono funzionalmente completi. È una
sintesi per ricetta: sempre possibile, raramente economica.

\subsection{Minimizzazione algebrica di $F$}

La forma canonica P di $F$ ha 5 clausole e 15 letterali. L'algebra del
Capitolo~\ref{ch:boole1} la riduce:
\begin{align*}
F &= \nott{A}BC+A\nott{B}\,\nott{C}+A\nott{B}C+AB\nott{C}+ABC\\
 &= \nott{A}BC+A\,(\nott{B}\,\nott{C}+\nott{B}C+B\nott{C}+BC) && \text{distributiva: si raccoglie $A$}\\
 &= \nott{A}BC+A\,\bigl(\nott{B}\,(\nott{C}+C)+B\,(\nott{C}+C)\bigr) && \text{distributiva, due volte}\\
 &= \nott{A}BC+A\,(\nott{B}+B) && \text{complemento: $\nott{C}+C=1$}\\
 &= \nott{A}BC+A && \text{complemento: $\nott{B}+B=1$}\\
 &= A+BC && \text{semplificazione: $x+\nott{x}\,y=x+y$.}
\end{align*}
Quindici letterali sono diventati tre: da cinque porte AND a tre ingressi e un
OR a cinque, a un AND e un OR. Funziona, ma ogni passo richiede un'idea. Ciò
che serve è un metodo.

\subsection{Lo scopo della minimizzazione}

Lo scopo è trovare, tra tutte le possibili forme a due livelli di tipo P, una
forma elementare \emph{minima}. Porte meno numerose e più piccole significano
costo inferiore, meno potenza e un circuito più veloce. Tre metodi sono
classici.

\begin{itemize}
\item \textbf{Mappe di Karnaugh}: grafico, per gruppi di celle adiacenti di
  area massima; pratico fino a quattro o cinque variabili
  \citep{karnaugh1953}.
\item \textbf{Metodo di Quine}: algebrico, per confronto sistematico dei termini.
\item \textbf{Metodo di McCluskey}: tabellare, adatto a un programma, e base
  degli odierni strumenti di sintesi logica. La procedura di Quine--McCluskey
  è esatta ma esponenziale; gli strumenti industriali usano euristiche.
\end{itemize}

% ----------------------------------------------------------------------------
\section{Mappe di Karnaugh}

\subsection{L'idea}

Una mappa di Karnaugh è la tabella di verità di una funzione di $n$ variabili
ridisegnata su una griglia di $2^{n}$ celle, una cella per riga. Le celle sono
ordinate in modo che due vicine differiscano in una sola variabile: di qui
l'ordine 00, 01, 11, 10 delle colonne, che è un codice Gray
(Capitolo~\ref{ch:combinational}), e di qui il fatto che in una mappa $m_3$
viene prima di $m_2$. Anche i bordi sinistro e destro sono adiacenti: la mappa
si richiude su sé stessa.

L'adiacenza grafica è così adiacenza logica: due 1 adiacenti sono due
mintermini che differiscono in una variabile, e per $x\,y+x\,\nott{y}=x$ si
fondono in una clausola con un letterale in meno.

\begin{prgconcept}[La regola della mappa]
Raggruppate gli 1 in rettangoli di 1, 2, 4 o 8 celle, i più grandi possibile,
finché ogni 1 è coperto. Un gruppo di $2^{k}$ celle elimina $k$ variabili: la
clausola di un gruppo contiene solo le variabili che al suo interno non
cambiano.
\end{prgconcept}

\subsection{Due, tre e quattro variabili}

\begin{figure}[htb]
\centering
\begin{subfigure}[b]{0.25\textwidth}\centering
\resizebox{0.8\textwidth}{!}{\begin{karnaugh-map}[2][2][1][$B$][$A$]
  \minterms{1,2,3}\autoterms[0]
  \implicant{2}{3}\implicant{1}{3}
\end{karnaugh-map}}
\caption{$F=A+B$}\label{fig:b2:k2}
\end{subfigure}\hfill
\begin{subfigure}[b]{0.36\textwidth}\centering
\resizebox{\textwidth}{!}{\begin{karnaugh-map}[4][2][1][$BC$][$A$]
  \minterms{3,4,5,6,7}\autoterms[0]
  \implicant{4}{6}\implicant{3}{7}
\end{karnaugh-map}}
\caption{$F=A+BC$}\label{fig:b2:k3}
\end{subfigure}\hfill
\begin{subfigure}[b]{0.36\textwidth}\centering
\resizebox{\textwidth}{!}{\begin{karnaugh-map}[4][4][1][$CD$][$AB$]
  \minterms{0,2,5,7,8,10,13,15}\autoterms[0]
  \implicant{5}{15}\implicantcorner
\end{karnaugh-map}}
\caption{$F=BD+\nott{B}\,\nott{D}$}\label{fig:b2:k4}
\end{subfigure}
\caption{Mappe di Karnaugh di due, tre e quattro variabili.}
\label{fig:b2:kmaps}
\end{figure}

\paragraph{Due variabili (Figura~\ref{fig:b2:k2}).} La funzione vale 1 per
$AB=01$, 10 e 11; la sua forma canonica $\nott{A}B+A\nott{B}+AB$ ha sei
letterali. Un gruppo è la riga $A=1$: al suo interno $B$ cambia, quindi $B$
sparisce e la clausola è $A$. L'altro è la colonna $B=1$: cambia $A$, e la
clausola è $B$. Dunque $F=A+B$: la funzione è semplicemente OR. La cella $AB$
appartiene a entrambi i gruppi; usare un 1 due volte è la legge di
idempotenza, $x+x=x$.

\paragraph{Tre variabili (Figura~\ref{fig:b2:k3}).} È la funzione $F$ del
capitolo: un 1 nelle celle dei mintermini 3, 4, 5, 6 e 7. Un gruppo di quattro
copre l'intera riga $A=1$: $B$ e $C$ cambiano entrambe, e la clausola è $A$.
Un gruppo di due copre la colonna $BC=11$: cambia $A$, e la clausola è $BC$.
Ogni 1 è coperto e nessun gruppo può essere ingrandito: $F=A+BC$, lo stesso
risultato dell'algebra, trovato a occhio in pochi secondi.

\paragraph{Quattro variabili (Figura~\ref{fig:b2:k4}).} La funzione
$F=\sum m(0,2,5,7,8,10,13,15)$ ha otto mintermini e trentadue letterali. Un
gruppo di quattro al centro ha dappertutto $B=1$ e $D=1$: la clausola è $BD$.
Anche i quattro angoli sono un solo gruppo, perché la mappa si richiude sia in
orizzontale sia in verticale, come se fosse disegnata su un toro: lì $B=0$ e
$D=0$, e la clausola è $\nott{B}\,\nott{D}$. Il risultato ha quattro letterali,
e si riconosce: $F$ è lo XNOR di $B$ e $D$.

La regola pratica è: prendere sempre prima i gruppi più grandi, poi coprire ciò
che resta.

\subsection{Condizioni di indifferenza}

A volte una combinazione degli ingressi non può mai presentarsi, oppure quando
si presenta l'uscita non importa. Una cifra in decimale codificato in binario,
per esempio, usa le configurazioni da 0000 a 1001, e le sei configurazioni da
1010 a 1111 non compaiono mai. Quelle celle si segnano nella mappa con un
trattino (o, in molti testi, con una croce): \term{condizione di indifferenza}
\eng{don't care}. Una cella di indifferenza può essere trattata come un 1
quando rende un gruppo più grande, e come uno 0 altrimenti. Più libertà
significa gruppi più grandi e un circuito meno costoso. L'esercizio classico è
il decodificatore per un display a sette segmenti, dove ogni segmento è una
funzione di quattro variabili con sei condizioni di indifferenza.

% ----------------------------------------------------------------------------
\section{Ritorno alla macchina}

\subsection{Trentadue sommatori completi fanno una \texttt{add}}

L'istruzione \code{add t2, t0, t1} porta i due registri a trentadue sommatori
completi, uno per posizione di bit (Figura~\ref{fig:b2:ripple}). Il riporto
uscente di ogni posizione è il riporto entrante della successiva: un
\term{sommatore a propagazione del riporto} \eng{ripple-carry adder}, in cui il
riporto si propaga verso sinistra come in una somma fatta a mano.

\begin{figure}[htb]
\centering
\begin{tikzpicture}[fa/.style={blkacc,minimum width=17mm,minimum height=12mm}]
  \foreach \i/\x/\l in {0/10.5/0, 1/7.9/1, 2/5.3/2, 3/1.3/31}{
    \node[fa] (f\i) at (\x,0) {sommatore\\\footnotesize bit \l};
    \draw[<-] ([xshift=-3mm]f\i.north) -- ++(0,6mm) node[above,font=\ttfamily\small]{A\l};
    \draw[<-] ([xshift=3mm]f\i.north) -- ++(0,6mm) node[above,font=\ttfamily\small]{B\l};
    \draw[->] (f\i.south) -- ++(0,-6mm) node[below,font=\ttfamily\small]{S\l};
  }
  \draw[->,very thick,prgamber] (12.4,0) node[right,font=\ttfamily]{0} -- (f0.east);
  \draw[->,very thick,prgamber] (f0.west) -- (f1.east);
  \draw[->,very thick,prgamber] (f1.west) -- (f2.east);
  \draw[->,very thick,prgamber] (f2.west) -- ++(-0.8,0);
  \node[text=prgamber,font=\bfseries] at (3.3,0) {$\cdots$};
  \draw[->,very thick,prgamber] (3.0,0) -- (f3.east);
  \draw[->,very thick,prgamber] (f3.west) -- ++(-0.9,0) node[left,lbl,text=prgamber,align=right]{riporto\\uscente};
  \node[lbl,text=prgamber] at (6.6,-0.9) {\emph{il riporto si propaga verso sinistra}};
\end{tikzpicture}
\caption{Un sommatore a propagazione del riporto di 32 bit.}\label{fig:b2:ripple}
\end{figure}

Il riporto uscente del bit 31 è l'overflow senza segno del
Capitolo~\ref{ch:integers}, che \riscv semplicemente scarta. I processori reali
usano schemi più veloci per il riporto, come il carry look-ahead
(Capitolo~\ref{ch:combinational}), ma la funzione calcolata è la stessa.

\subsection{Un sommatore completo, una porta per istruzione}

\asmfile[caption={Il sommatore completo come programma.},label=lst:b2:fa]{code/l09/s27.S}

Cinque istruzioni, cinque porte: il Listato~\ref{lst:b2:fa} è lo schema della
Figura~\ref{fig:b2:fa}, letto da sinistra a destra. Eseguirlo passo-passo in
\prgqt e compilare la tabella di verità è \emph{analisi}; ricavare queste porte
dalla tabella con il metodo di questo capitolo era \emph{sintesi}. E con
registri interi come ingressi, le stesse cinque righe lavorano su tutte le
trentadue posizioni di bit insieme.

% ----------------------------------------------------------------------------
\section{Verificate di aver capito}

\begin{quiz}{Quante funzioni booleane distinte di tre variabili esistono?}
\item 8
\item 9
\item 64
\item 256
\end{quiz}

\begin{quiz}{Quale mintermine di tre variabili $A$, $B$, $C$ corrisponde alla riga 5 della tabella di verità?}
\item $\nott{A}\,B\,\nott{C}$
\item $A\,\nott{B}\,C$
\item $A+\nott{B}+C$
\item $A\,B\,\nott{C}$
\end{quiz}

\begin{quiz}{In una mappa di Karnaugh, un gruppo rettangolare di quattro 1 adiacenti elimina:}
\item una variabile
\item due variabili
\item tre variabili
\item quattro variabili
\end{quiz}

\begin{quiz}{In un sommatore completo, il riporto uscente vale 1 quando:}
\item esattamente un ingresso vale 1
\item un numero dispari di ingressi vale 1
\item almeno due ingressi valgono 1
\item solo quando tutti e tre gli ingressi valgono 1
\end{quiz}

\section{Esercizi}

\begin{exercise}[prima su carta, poi su \prg]
\begin{enumerate}[nosep]
\item Scrivete le forme canoniche P e S delle uscite $S$ e $C_{out}$ del
  sommatore completo.
\item Minimizzate $C_{out}$ con una mappa di Karnaugh di tre variabili. Quante
  porte si risparmiano?
\item Minimizzate $F=\sum m(0,1,2,5,8,9,10)$ con una mappa di quattro variabili.
\item Costruite NOT, AND e OR usando solo NAND, e verificate ciascuno con una
  tabella di verità.
\item Programmate il sommatore completo su \prg e mostrate $S$ e $C_{out}$
  come lampade per le otto combinazioni degli ingressi.
\end{enumerate}
\expected{Su carta: la forma canonica di $C_{out}$ ha quattro mintermini (3, 5,
6, 7) e la mappa dà tre clausole di due letterali,
$C_{out}=AB+AC_{in}+BC_{in}$, la funzione di \emph{maggioranza}. Il terzo
compito dà $F=\nott{B}\,\nott{D}+\nott{B}\,\nott{C}+\nott{A}\,\nott{C}\,D$. Su
\prg: otto coppie di lampade che riproducono la tabella di verità del sommatore
completo.}
\end{exercise}

\section{Glossario}

\begin{glossarytable}
literal & letterale\\
clause, product term & clausola, termine prodotto\\
factor, sum term & fattore, termine somma\\
minterm / maxterm & mintermine / maxtermine\\
half adder / full adder & semisommatore / sommatore completo\\
carry in / carry out & riporto entrante / uscente\\
canonical form & forma canonica\\
sum of products (SOP) & somma di prodotti (tipo P)\\
product of sums (POS) & prodotto di somme (tipo S)\\
Karnaugh map & mappa di Karnaugh\\
don't care & condizione di indifferenza\\
ripple-carry adder & sommatore a propagazione del riporto\\
\end{glossarytable}

\section{Conclusioni e letture}

Una funzione booleana è una tabella di verità, e una tabella di verità dà
sempre un'espressione canonica. La minimizzazione è il processo che ottimizza
una rete logica per semplificarne la realizzazione. Le mappe di Karnaugh
trasformano l'algebra in un disegno: raggruppare gli 1, leggere le clausole. Il
semisommatore e il sommatore completo sono i primi veri circuiti di questo
libro, e trentadue sommatori completi in fila sono ciò che esegue la \ins{add}
scritta nel Capitolo~\ref{ch:firstcontact}.

\paragraph{Letture.} \citet{patterson2020}, appendice~A; \citet{bucci2017}, i
capitoli sulle reti combinatorie; Charles Petzold, \emph{Code}, i capitoli
sulle porte logiche e sulle macchine addizionatrici binarie.
